%%% List of acronyms and abbreviations
% To define a new acronym:
% \newabbreviation{<key>}{<abbreviation>}{<long form>}
%
% To make it easier to use, also define a new command with the format:
% \newcommand{\<key>}{\gls*{<key>}}
%
% The * in \gls*{<key>} is to prevent LaTeX from adding hyperlinks if the
% hyperref package is used.
%
% In the text, just use \<key>{}. On the first usage, LaTeX will display the
% long form followed by the acronym in parenthesis. Then, it will only display
% the acronym.
%
% For the special case where you want the long version later in the text, use
% the \gls family of commands.
%%%
% separate the abbreviations from the acronyms
\usepackage[automake,nonumberlist,abbreviations,acronyms]{glossaries-extra}
\makeglossaries
\usepackage{glossaries-extra}

\newabbreviation{ci}{CI}{Contextual Integrity}
\newcommand{\ci}{\gls*{ci}\@}
\newabbreviation{fipps}{FIPPs}{Fair Information Practice Principles}
\newcommand{\fipps}{\gls*{fipps}\@}
\newabbreviation{ftc}{FTC}{U.S. Federal Trade Commission}
\newcommand{\ftc}{\gls*{ftc}\@}
\newabbreviation{gdpr}{GDPR}{General Data Protection Regulation}
\newcommand{\gdpr}{\gls*{gdpr}\@}
\newabbreviation{pii}{PII}{Personally Identifiable Information}
\newcommand{\pii}{\gls*{pii}\@}
\newabbreviation{pipeda}{PIPEDA}{the Personal Information Protection and Electronic Documents Act}
\newcommand{\pipeda}{\gls*{pipeda}\@}
\newabbreviation{sdk}{SDK}{Software Development Kit}
\newcommand{\sdk}{\gls*{sdk}\@}
\newcommand{\sdks}{\glspl*{sdk}\@}

